\section{章节标题}

\subsection{本节标题}

% 本节问题：这一小节要回答的唯一问题/判断是什么？
% 输入/证据：使用哪些题面条件、数据、结果文件、图表或文献？
% 方法/定义：需要定义哪些变量、公式、算法、比较口径或约束？
% 结果：本节必须交付哪个数字、规律、方案或负结果？
% 解释/边界：为什么出现该结果，在什么条件下成立，代价是什么？
% 接口：本节输出是否进入下一节/下一问；没有则写“无”。
在此直接撰写可编译的 LaTeX 正文。公式使用数学环境，图表使用唯一的语义化 \texttt{label}。结果段按“回答的问题—数字或图表—同口径比较—原因解释—适用边界/代价”组织。

% Figure 叙述模板（删除方括号提示后再提交）：
% 图前定位：图~\ref{fig:semantic-name} 回答【唯一问题】，比较口径为【对象/范围/单位】，重点看【面板/颜色/线型】。
% \begin{figure}[htbp]
%   \centering
%   \includegraphics[width=\linewidth]{figures/正式高清图}
%   \caption{【对象与范围】；【面板/颜色/线型/箭头语义】；【必要统计口径】。}
%   \label{fig:semantic-name}
% \end{figure}
% 图后解释：从图中读出【具体数值/趋势/关系/异常】，说明【对本节结论的意义】；但该图不能证明【边界】，因此下一步【验证/转折】。
